\documentclass[
  aps,
  prd,
  reprint,
  superscriptaddress,
  floatfix,
  nofootinbib,
  longbibliography
]{revtex4-2}

\usepackage[T1]{fontenc}
\usepackage[utf8]{inputenc}
\usepackage{amsmath,amssymb,amsfonts,amsthm,mathtools}
\usepackage{microtype}
\usepackage{hyperref}

\hypersetup{
  colorlinks=true,
  linkcolor=blue,
  citecolor=blue,
  urlcolor=blue
}

\theoremstyle{plain}
\newtheorem{theorem}{Theorem}
\newtheorem{proposition}[theorem]{Proposition}
\newtheorem{corollary}[theorem]{Corollary}
\newtheorem{criterion}[theorem]{Criterion}
\theoremstyle{definition}
\newtheorem{definition}[theorem]{Definition}
\theoremstyle{remark}
\newtheorem{remark}[theorem]{Remark}

\newcommand{\dd}{\mathrm d}
\newcommand{\D}{\mathcal D}
\newcommand{\vis}{\mathrm{vis}}
\newcommand{\wind}{\mathrm w}
\newcommand{\eff}{\mathrm{eff}}
\newcommand{\crit}{\mathrm{crit}}
\newcommand{\GI}{\mathrm{GI}}
\newcommand{\Ok}{\Omega_K}
\newcommand{\Ow}{\Omega_{\mathrm w}}
\newcommand{\Oq}{\Omega_q}
\newcommand{\Aq}{\mathcal A_q}
\newcommand{\Aw}{\mathcal A_{\mathrm w}}
\newcommand{\Rcl}{\mathcal R}
\DeclareMathOperator{\Vol}{Vol}
\DeclareMathOperator{\Spec}{Spec}
\DeclareMathOperator{\Tr}{Tr}
\newcommand{\id}{\mathbf 1}
\newcommand{\Z}{\mathbb Z}
\newcommand{\R}{\mathbb R}
\newcommand{\eps}{\epsilon}
\newcommand{\qhat}{\widehat q}
\newcommand{\NF}{N_F}
\newcommand{\eF}{e_F}
\newcommand{\PG}{P_\Gamma}
\newcommand{\Hn}{\mathcal H_n}
\newcommand{\LCDM}{\Lambda\mathrm{CDM}}

\begin{document}

\title{Compact Spherical FLRW Residuals: Closure, Local \texorpdfstring{$\Lambda$}{Lambda}CDM Degeneracy, and Amplitude-Selection Boundaries}

\author{Andrei~Preece}
\email{andre.preece@web.de}
\affiliation{Independent Researcher, Berlin, Germany}

\author{Boris~Batenin}
\email{batenin@web.de}
\affiliation{Independent Researcher, Berlin, Germany}

\date{May 15, 2026}

\begin{abstract}
For a fixed spacetime metric and a declared visible-sector subtraction, the
Einstein residual
\[
T^\D_{\mu\nu}=\frac{c^4}{8\pi G}G_{\mu\nu}-T^{\vis}_{\mu\nu}
\]
is covariantly fixed.  We prove that, within the homogeneous locally-FLRW
residual class, this residual has the minimal pressureless-current plus
fixed-four-form stress-tensor form, at the homogeneous stress-tensor level,
precisely when
\[
\frac{\dd}{\dd\ln a}[-p_\D]=0,
\qquad
\frac{\dd}{\dd\ln a}\left\{a^3(\eps_\D+p_\D)\right\}=0,
\]
with the corresponding constants nonnegative.  Equivalently,
\[
\eps_\D(a)=\eps_{\wind,0}a^{-3}+\eps_q,
\qquad
p_\D(a)=-\eps_q .
\]
Thus the minimal homogeneous branch is locally identical to a cold-dark-matter
plus cosmological-constant residual with the same amplitudes and curvature.
Compact spherical topology may still affect global volume, quotient-projected
spectra, compact covariance, and observer-position
dependence for non-globally-homogeneous quotients, but it does not by itself determine \(\Omega_c\),
\(\Omega_\Lambda\), or \(\Lambda_{\eff}\).  The result is therefore a
residual-closure theorem and an amplitude-selection boundary, not an
abundance-origin theorem or a topology detection.
\end{abstract}

\maketitle

\section{Claim boundary}

This paper proves a conditional stress-tensor closure theorem.  It does not
derive the observed dark-sector abundances, does not select a flux branch, and
does not claim detection of compact topology.  Once the metric
\(g_{\mu\nu}\), the Einstein-side convention, and the visible-sector subtraction
\(T^{\vis}_{\mu\nu}\) are fixed, the Einstein residual is fixed.  If this
residual is homogeneous and locally FLRW, then it has the minimal
pressureless-current plus fixed-vacuum stress-tensor form if and only if
\begin{equation}
-p_\D=\text{constant}\ge0,
\qquad
a^3(\eps_\D+p_\D)=\text{constant}\ge0 .
\label{eq:claim_boundary_closure}
\end{equation}
Microscopic current and four-form realizations may impose further amplitude
selection rules; the closure theorem itself does not supply them.  The
corresponding stress tensor is locally indistinguishable from cold dark matter
plus \(\Lambda\) with the same amplitudes and curvature.  Any distinction must
therefore come from compact spectral projection, quotient covariance, perturbative null tests,
or an additional amplitude-selection rule.

The paper is complete only at this declared level: homogeneous residual
closure, compact-background bookkeeping, perturbative null conditions, and
amplitude-selection boundaries.  It is not a compact Boltzmann solver, a
likelihood analysis, or a microscopic dark-sector completion.

\section{Einstein residual and convention covariance}

Let \((M,g_{\mu\nu})\) be a four-dimensional spacetime, and let
\(T^{\vis}_{\mu\nu}\) denote the stress tensor assigned to the declared visible
sector.  Define
\begin{equation}
\boxed{
T^\D_{\mu\nu}:=
\frac{c^4}{8\pi G}G_{\mu\nu}-T^{\vis}_{\mu\nu} .
}
\label{eq:primitive_residual_v5}
\end{equation}
A bare cosmological-constant term is included in \(T^\D_{\mu\nu}\) unless the
Einstein-side convention is explicitly changed to put \(G_{\mu\nu}+\Lambda
 g_{\mu\nu}\) on the left-hand side.

\begin{proposition}[Residual uniqueness and convention covariance]
If
\begin{equation}
G_{\mu\nu}=\frac{8\pi G}{c^4}\left(T^{\vis}_{\mu\nu}+S_{\mu\nu}\right),
\end{equation}
then \(S_{\mu\nu}=T^\D_{\mu\nu}\).  If the visible subtraction is changed by
\(X_{\mu\nu}\), then
\begin{equation}
T^{\vis}_{\mu\nu}\to T^{\vis}_{\mu\nu}+X_{\mu\nu},
\qquad
T^\D_{\mu\nu}\to T^\D_{\mu\nu}-X_{\mu\nu} .
\end{equation}
If the Einstein side is changed from \(G_{\mu\nu}\) to
\(G_{\mu\nu}+\Lambda_0g_{\mu\nu}\) while the visible subtraction is kept fixed,
then
\begin{equation}
T^\D_{\mu\nu}\to T^\D_{\mu\nu}+\frac{c^4\Lambda_0}{8\pi G}g_{\mu\nu} .
\end{equation}
\end{proposition}

We use standard differential-geometric conventions~\cite{Wald}.  The contracted Bianchi identity gives
\begin{equation}
\nabla^\mu T^\D_{\mu\nu}=-\nabla^\mu T^{\vis}_{\mu\nu} .
\label{eq:bianchi_residual_v5}
\end{equation}
Thus separate residual conservation is not automatic; it follows only when the
visible sector is separately conserved in the adopted convention.  This point
matters because the minimal dust-plus-vacuum source branch is separately
conserved.

We use a timelike four-velocity \(U^\mu\) normalized by
\(U^\mu U_\mu=-c^2\) and the projector
\begin{equation}
h_{\mu\nu}=g_{\mu\nu}+\frac{U_\mu U_\nu}{c^2} .
\label{eq:projector_v5}
\end{equation}
For a perfect fluid with physical energy density \(\eps\) and pressure \(p\),
\begin{equation}
T_{\mu\nu}=\eps\frac{U_\mu U_\nu}{c^2}+p h_{\mu\nu} .
\label{eq:perfect_fluid_v5}
\end{equation}

\section{Locally-FLRW reconstruction}

On the round cover, use
\begin{equation}
\dd s^2=-c^2\dd t^2+a^2(t)A_0^2\dd\Omega_3^2,
\qquad
\Ok=-\frac{c^2}{H_0^2A_0^2}<0,
\label{eq:closed_metric_v5}
\end{equation}
where \(A_0\) is the present curvature radius of the simply connected cover.
For a quotient \(S^3_{A_0}/\Gamma\), the local metric is unchanged while the
volume and global spectrum are restricted.  Let
\begin{equation}
E(a)=\frac{H(a)}{H_0},
\qquad
\eps_{\crit,0}=\frac{3H_0^2c^2}{8\pi G} .
\end{equation}
The local Friedmann equations reconstruct a homogeneous residual as
\begin{align}
\eps_\D(a)
&=\eps_{\crit,0}\left[E^2(a)-\Ok a^{-2}\right]
-\eps_{\vis}(a),
\label{eq:eps_reconstruction_v5}\\
p_\D(a)
&=-\frac{\eps_{\crit,0}}{3}\biggl[
2E(a)\frac{\dd E}{\dd\ln a}
\nonumber\\
&\hspace{1.0cm}
+3E^2(a)-\Ok a^{-2}
\biggr]-p_{\vis}(a).
\label{eq:p_reconstruction_v5}
\end{align}
These are reconstruction identities, not equations of microscopic origin.

\begin{theorem}[Locally-FLRW residual reduction]
After a fixed visible subtraction, every homogeneous and isotropic symmetric
residual tensor is uniquely of the form
\begin{equation}
T^\D_{\mu\nu}=\eps_\D(a)\frac{U_\mu U_\nu}{c^2}+p_\D(a)h_{\mu\nu} .
\label{eq:flrw_residual_tensor_v5}
\end{equation}
If the residual is separately conserved, then
\begin{equation}
\frac{\dd\eps_\D}{\dd\ln a}+3(\eps_\D+p_\D)=0 .
\label{eq:residual_conservation_v5}
\end{equation}
\end{theorem}

\begin{proof}
Local homogeneity and isotropy leave only \(U_\mu U_\nu\) and the spatial
projector \(h_{\mu\nu}\) as invariant symmetric rank-two tensors.  Taking the
covariant divergence in FLRW gives Eq.~\eqref{eq:residual_conservation_v5}.
\end{proof}

\section{Closure-pair atlas and minimal branch}

Define the closure variables
\begin{equation}
\boxed{
\Aq(a):=-p_\D(a),
\qquad
\Aw(a):=a^3[\eps_\D(a)+p_\D(a)] .
}
\label{eq:closure_pair_v5}
\end{equation}
They are a lossless reparametrization of the homogeneous residual:
\begin{equation}
p_\D(a)=-\Aq(a),
\qquad
\eps_\D(a)=a^{-3}\Aw(a)+\Aq(a) .
\label{eq:atlas_inverse_v5}
\end{equation}
If the residual is separately conserved, then
\begin{equation}
\frac{\dd\Aq}{\dd\ln a}+a^{-3}\frac{\dd\Aw}{\dd\ln a}=0 .
\label{eq:atlas_flow_v5}
\end{equation}
Thus the general separately conserved homogeneous residual is a one-function
flow in \((\Aq,\Aw)\) space.  The minimal dust-plus-vacuum branch is the
constant nonnegative submanifold.

\begin{theorem}[Minimal homogeneous residual closure]
\label{thm:minimal_closure_v5}
Let \(I\subset\R_+\) be a connected scale-factor interval, and let
\((\eps_\D,p_\D)\in C^1(I)\) be the homogeneous locally-FLRW Einstein residual
after a fixed visible-sector subtraction.  Then the residual is representable
at the homogeneous stress-tensor level as
\begin{equation}
T^\D_{\mu\nu}
=
\eps_{\wind,0}a^{-3}\frac{U_\mu U_\nu}{c^2}
-
\eps_q g_{\mu\nu},
\qquad
\eps_{\wind,0},\eps_q\ge0,
\label{eq:minimal_tensor_v5}
\end{equation}
if and only if \(\Aq\) and \(\Aw\) are constant and nonnegative on \(I\).  In
that case
\begin{equation}
\begin{gathered}
\eps_q=\Aq,\qquad \eps_{\wind,0}=\Aw,\\
\eps_\D(a)=\eps_{\wind,0}a^{-3}+\eps_q,
\qquad
p_\D(a)=-\eps_q .
\end{gathered}
\label{eq:closure_solution_v5}
\end{equation}
\end{theorem}

\begin{proof}
If Eq.~\eqref{eq:minimal_tensor_v5} holds, then
\(-p_\D=\eps_q\) and \(a^3(\eps_\D+p_\D)=\eps_{\wind,0}\).  Conversely, if
\(\Aq\) and \(\Aw\) are constant, then Eq.~\eqref{eq:atlas_inverse_v5} gives
\(p_\D=-\Aq\) and \(\eps_\D=\Aw a^{-3}+\Aq\), which is precisely
Eq.~\eqref{eq:minimal_tensor_v5}.  Nonnegativity is exactly the condition
\(\eps_q,\eps_{\wind,0}\ge0\).
\end{proof}

The ``if and only if'' is an iff statement for the homogeneous stress tensor.
A microscopic branch realizes the constants only if its charge, mass scale,
flux normalization, branch labels, and quotient volume admit those amplitudes.
In particular, a discrete flux lattice can make some otherwise admissible
constant values of \(\eps_q\) unavailable unless the flux unit or branch
spectrum is correspondingly specified.

The invariant null tests are
\begin{equation}
\begin{gathered}
\Rcl_q:=\frac{\dd\Aq}{\dd\ln a}=0,
\qquad
\Rcl_{\wind}:=\frac{\dd\Aw}{\dd\ln a}=0,\\
\Aq,\Aw\ge0 .
\end{gathered}
\label{eq:null_tests_v5}
\end{equation}
A failure of these tests falsifies the minimal homogeneous branch under the
adopted visible subtraction; it does not falsify compact topology.  Since the
closed minimal branch is separately conserved, Eq.~\eqref{eq:bianchi_residual_v5}
also implies separate conservation of the declared visible sector in the same
convention.  Failure of visible-sector conservation is therefore already a
failure of the adopted minimal-source decomposition, not an independent
topology effect.

\begin{corollary}[Local \(\LCDM\) degeneracy]
On any local FLRW chart, the minimal branch is exactly the stress tensor of a
pressureless component plus a cosmological-constant component with the same
amplitudes and curvature.  No homogeneous observable depending only on
\(H(a)\), distances, and the same curvature can distinguish the minimal branch
from CDM plus \(\Lambda\).  A CDM plus \(\Lambda\) model placed on the same
compact quotient, with the same projected perturbation covariance, the same
initial data, and the same amplitudes, remains degenerate with the minimal
branch at this level.
\end{corollary}

\section{Current and four-form representation}

The minimal source representation consists of a pressureless conserved-current
branch and a fixed four-form branch.  The current branch may carry a protected
degree label, but topology alone does not give a dust stress tensor.

The degree label is standard in algebraic topology~\cite{Hatcher}, while the
pressureless current representative is a Brown-fluid effective
description~\cite{BrownFluid}.  Let \(\Sigma\) be compact and oriented, let
\(\phi:\Sigma\to S^3\), and normalize the target volume form by
\(\int_{S^3}\widehat\Omega=1\).  The degree is
\begin{equation}
Q=\int_\Sigma \phi^*\widehat\Omega\in\Z .
\label{eq:degree_v5}
\end{equation}
The associated topological current is identically conserved, but
\begin{equation}
Q\in\Z\not\Rightarrow N^\mu=n(a)u^\mu\not\Rightarrow p_{\wind}=0 .
\label{eq:q_not_dust_v5}
\end{equation}
The dust result requires a homogeneous coarse-grained Brown-current
representative and a linear energy density
\begin{equation}
\eps_{\wind}(n)=\mu_{\wind}c^2n .
\label{eq:linear_current_v5}
\end{equation}
More generally, a barotropic current effective theory with
\(\eps_J=F(n)\) has
\begin{equation}
p_J=nF'(n)-F(n),
\qquad
\frac{v_{s,J}^2}{c^2}:=\frac{\dd p_J}{\dd\eps_J}
=\frac{nF''(n)}{F'(n)} .
\label{eq:current_atlas_v5}
\end{equation}
The minimal pressureless branch is the special case \(F(n)=\mu_{\wind}c^2n\),
for which \(p_J=0\) and \(v_{s,J}=0\).

For a homogeneous coarse-grained current branch on
\(\Sigma_{A_0}=S^3_{A_0}/\Gamma\),
\begin{equation}
\Vol(\Sigma_{A_0})=\frac{2\pi^2A_0^3}{|\Gamma|},
\qquad
n(a)=\frac{|Q|}{a^3\Vol(\Sigma_{A_0})},
\label{eq:volume_number_density_v5}
\end{equation}
and therefore
\begin{equation}
\Ow=
\frac{\mu_{\wind}c^2|Q|}{\Vol(\Sigma_{A_0})\eps_{\crit,0}}
=
\frac{4G\mu_{\wind}|Q|\,|\Gamma|}{3\pi H_0^2A_0^3} .
\label{eq:omega_w_v5}
\end{equation}
This is an exact bookkeeping identity once \((Q,\mu_{\wind},|\Gamma|,A_0)\)
are supplied.  It is not an amplitude prediction.

Four-form vacuum sectors of this kind are standard in gravitational theory~\cite{Aurilia,HenneauxTeitelboim,BrownTeitelboim,BoussoPolchinski}.  For a four-form \(F_4=\dd C_3\), take
\begin{equation}
S_F=-\frac{Z_F}{2\cdot4!}\int \sqrt{-g}\,
F_{\mu\nu\rho\sigma}F^{\mu\nu\rho\sigma}\,\dd^4x,
\qquad Z_F>0 .
\label{eq:four_form_action_v5}
\end{equation}
The field equation implies that the four-form is proportional to the metric
volume form on each connected branch,
\begin{equation}
F_{\mu\nu\rho\sigma}=f\varepsilon_{\mu\nu\rho\sigma},
\qquad
f=\text{constant},
\label{eq:four_form_branch_v8}
\end{equation}
up to the orientation convention for the Hodge dual.  Only \(f^2\) enters the
stress tensor.  Define
\begin{equation}
\qhat=\sqrt{Z_F}f,
\qquad
\eps_q=\frac12\qhat^2 .
\label{eq:four_form_energy_v5}
\end{equation}
The stress tensor obtained from Eq.~\eqref{eq:four_form_action_v5} is
\begin{equation}
T^F_{\mu\nu}
=
\frac{Z_F}{3!}F_{\mu\alpha\beta\gamma}
F_{\nu}{}^{\alpha\beta\gamma}
-
\frac{Z_F}{2\cdot4!}g_{\mu\nu}F_{\alpha\beta\gamma\delta}
F^{\alpha\beta\gamma\delta} .
\label{eq:four_form_stress_general_v7}
\end{equation}
With the \((-+++)\) signature,
\begin{equation}
\varepsilon_{\mu\alpha\beta\gamma}
\varepsilon_{\nu}{}^{\alpha\beta\gamma}=-3!g_{\mu\nu},
\qquad
F_{\alpha\beta\gamma\delta}F^{\alpha\beta\gamma\delta}=-4!f^2,
\end{equation}
which gives
\begin{equation}
T^q_{\mu\nu}=-\eps_q g_{\mu\nu},
\qquad
\Lambda_{\eff}=\frac{8\pi G}{c^4}\eps_q .
\label{eq:four_form_stress_v5}
\end{equation}
For a fixed superselection branch, the mixed-index perturbation vanishes,
\begin{equation}
\delta T^{q\,\mu}{}_{\nu}=0,
\label{eq:mixed_index_four_form_v5}
\end{equation}
while the covariant-index variation
\begin{equation}
\delta T^q_{\mu\nu}=-\eps_q\delta g_{\mu\nu}
\label{eq:covariant_four_form_v5}
\end{equation}
is purely metric-induced.

\begin{proposition}[Closure versus flux-lattice selection]
The homogeneous closure criterion determines only the stress-energy form
\begin{equation}
T^q_{\mu\nu}=-\eps_q g_{\mu\nu},
\qquad
\eps_q=\text{constant}\ge0 .
\end{equation}
It does not by itself determine whether the constant belongs to a discrete flux
lattice.  If a compact zero-mode normalization is supplied so that
\begin{equation}
\qhat=\NF\eF,
\qquad
\NF\in\Z,
\end{equation}
then the admissible vacuum amplitudes are restricted to
\begin{equation}
\eps_q=\frac12\NF^2\eF^2 .
\label{eq:flux_lattice_v5}
\end{equation}
Thus the FLRW closure theorem is a continuous stress-energy classification,
whereas flux quantization is an additional amplitude-selection constraint.
Arbitrary residual matching requires an unfixed normalization \(\eF\), a dense
landscape of allowed branches, or a phenomenological effective constant.
\end{proposition}

\section{Compact topology and quotient projection}

The spatial topology is a global input, not a local equation of state.
Throughout the topological-current and compact three-form slots we restrict to
compact orientable spherical space forms, equivalently to finite freely acting
orientation-preserving isometry groups \(\Gamma\) on \(S^3\).  Without
orientability the degree and top-cohomology statements must be replaced by their
twisted-coefficient versions; such branches are not considered here.  The
classification of complete constant-curvature space forms gives~\cite{HawkingEllis,Wolf}
\begin{equation}
\Sigma=\widetilde\Sigma/\Gamma,
\qquad
\widetilde\Sigma\in\{S^3,\R^3,H^3\} .
\end{equation}
In the compact positive-curvature case,
\begin{equation}
\Sigma=S_A^3/\Gamma,
\qquad
\Vol(\Sigma)=\frac{2\pi^2A^3}{|\Gamma|} .
\label{eq:space_form_volume_v5}
\end{equation}
If simple connectivity is assumed, \(\Gamma=\id\) and \(\Sigma\simeq S_A^3\).
Without simple connectivity, the invariant conclusion is only that the
universal cover is \(S_A^3\).

On the simply connected cover,
\begin{equation}
-\Delta_{S_A^3}Y_{n\ell m}=\frac{n(n+2)}{A^2}Y_{n\ell m},
\qquad
{\rm mult}(n)=(n+1)^2 .
\label{eq:s3_spectrum_v5}
\end{equation}
For a quotient, fields are \(\Gamma\)-invariant cover fields.  The projector is
\begin{equation}
\boxed{
\PG=\frac1{|\Gamma|}\sum_{\gamma\in\Gamma}\gamma^* .
}
\label{eq:quotient_projector_v5}
\end{equation}
The physical scalar eigenspace is \(\Hn^\Gamma=\PG\Hn\), with multiplicity
\begin{equation}
m_\Gamma(n)=\Tr_{\Hn}\PG .
\label{eq:quotient_multiplicity_v5}
\end{equation}
The same logic applies to physical vector and tensor harmonic sectors after
closed-FLRW gauge and constraint modes are removed.

Compact topology can therefore enter through
\begin{equation}
\Vol(\Sigma),\qquad \Spec_\Gamma(\Delta),\qquad m_\Gamma(n),\qquad
\mathrm{Cov}_\Gamma,
\label{eq:topology_channels_v5}
\end{equation}
not through a new local homogeneous stress tensor.  Observer-position dependence
appears for non-globally-homogeneous quotients; for globally homogeneous
quotients the covariance is still quotient-projected but is not position
dependent in the same way.

\begin{criterion}[Topology-evidence boundary]
A compact spherical model may claim compact-background consistency after
specifying \(S^3/\Gamma\) and the quotient projection.  It may not claim topology
evidence or benchmark superiority until the projector has been propagated into
transfer functions and observable covariance,
\begin{equation}
\PG\longrightarrow
m_\Gamma(n),\ \Xi_n^\Gamma,\
\Delta_{\ell n}^{X,\Gamma},\ C_\ell^{XY,\Gamma},\
\mathrm{Cov}_\Gamma,\ \mathcal L_\Gamma .
\label{eq:transfer_boundary_v5}
\end{equation}
\end{criterion}

\section{Amplitude-selection boundary}

Compact topology supplies labels, volume, and spectral restrictions.  It does
not supply numerical cosmological abundances.  In particular,
\begin{equation}
(\Gamma,[\Sigma,S^3],H^3(\Sigma,\Z),Q,\NF)
\not\Rightarrow
(\Ow,\Oq,\Lambda_{\eff}) .
\label{eq:not_amplitudes_v5}
\end{equation}
Even after the geometric radius \(A_0\) is supplied, the amplitudes require
\(\mu_{\wind}\), \(\eF\), branch labels, compact-gauge normalization, and a
selection rule or phenomenological prior.  This is an amplitude-selection
boundary, not a failure of the closure theorem.

A useful diagnostic for selector strength is the curvature-linear ansatz
\(\Ow=C|\Omega_K|\).  If observations impose \(|\Omega_K|<\kappa\) and the
target is \(\Omega_\star\), then
\begin{equation}
C>\frac{\Omega_\star}{\kappa} .
\label{eq:selector_bound_v5}
\end{equation}
For \(\Omega_\star\simeq0.26\) and \(\kappa=0.005\), one needs \(C>52\).
A toy holographic winding saturation condition,
\begin{equation}
\mu_{\wind}c^2|Q|=\frac{c^4A_0}{2G},
\end{equation}
gives only
\begin{equation}
\Ow^{\rm HWS}=\frac{2}{3\pi}|\Omega_K|
\label{eq:hws_simple_v5}
\end{equation}
for the simply connected case.  For a quotient, imposing the same total-energy
saturation on the quotient volume gives
\begin{equation}
\Ow^{\rm HWS}(\Gamma)=\frac{2|\Gamma|}{3\pi}|\Omega_K| .
\label{eq:hws_quotient_v5}
\end{equation}
This is not a derivation of \(\Omega_c\).  If \(|\Gamma|\sim 10^2\) is used to
obtain a large effective coefficient, the proposal becomes a large-quotient
amplitude ansatz that must pass matched-circle, observer-position, covariance,
and spectral-multiplicity tests.

\section{Perturbative and observational null tests}

We use standard gauge-invariant cosmological perturbation conventions~\cite{Bardeen,KodamaSasaki,Mukhanov,MaBertschinger}.  At linear order, the residual is defined before gauge choice by
\begin{equation}
\delta T^\D_{\mu\nu}=\frac{c^4}{8\pi G}\delta G_{\mu\nu}
-\delta T^{\vis}_{\mu\nu} .
\label{eq:linear_residual_v5}
\end{equation}
The fixed four-form branch contributes no independent scalar perturbation on a
fixed branch.  After subtracting this fixed vacuum component, the minimal
current branch must obey the dust null conditions on every physical
quotient-projected scalar mode:
\begin{equation}
\boxed{
\delta p_{\wind,{\rm rf}}=0,
\qquad
\Pi_{\wind}^{\GI}=0,
\qquad
\frac{v_{s,\wind}^2}{c^2}=0 .
}
\label{eq:dust_null_conditions_v5}
\end{equation}
A statistically significant intrinsic pressure perturbation, rest-frame sound
speed, anisotropic stress, entropy source, or independent non-dust stress source
not assigned to the visible sector falsifies the minimal perturbative branch.

Current background data sharpen the closure tests without changing their
logical status.  Planck 2018 combined with BAO gives a curvature constraint
consistent with flatness, conventionally quoted as
\(\Omega_K=0.0007\pm0.0019\) for TT,TE,EE+lowE+lensing+BAO~\cite{Planck2018}.
Recent DESI DR2 BAO results are well described by flat \(\Lambda\)CDM, while
BAO+CMB prefers a \(w_0w_a\) model over \(\Lambda\)CDM at \(3.1\sigma\), and
BAO+CMB plus supernova samples gives a sample-dependent preference in the range
\(2.8\)--\(4.2\sigma\)~\cite{DESIDR2}.  If a reconstructed residual
pressure evolution remains significant after systematics, dataset tension, and
visible-sector convention dependence are accounted for, then the fixed-four-form
minimal branch fails the
\(\Rcl_q=0\) null test.  That would falsify the fixed-vacuum source law,
not compact spherical topology.

Likewise, matched-circle and compact-covariance limits constrain detectable
domains and topology families, not every compact spherical quotient.
Lens-space constraints are therefore constraints on specified \(S^3/\Gamma\)
implementations and their observer-position/covariance pipelines, where
applicable~\cite{Cornish,PlanckTopology,SahaLens}.

\section{Conclusion}

We have proved a closure theorem for homogeneous Einstein residuals on compact
spherical FLRW backgrounds.  After the metric, Einstein-side convention, and
visible-sector subtraction are fixed, the residual belongs to the minimal
pressureless-current plus fixed-four-form branch exactly when
\begin{equation}
-p_\D=\text{constant}\ge0,
\qquad
a^3(\eps_\D+p_\D)=\text{constant}\ge0 .
\end{equation}
Equivalently,
\begin{equation}
\eps_\D(a)=\eps_{\wind,0}a^{-3}+\eps_q,
\qquad
p_\D(a)=-\eps_q .
\end{equation}
This proves local background equivalence to cold dark matter plus \(\Lambda\)
with the same curvature and amplitudes.

Compact topology can still enter through finite volume, quotient-projected
spectra, compact covariance, and observer-position effects for
non-globally-homogeneous quotients.  It does
not, however, determine \(\Omega_c\), \(\Omega_\Lambda\), or
\(\Lambda_{\eff}\).  The Brown-current dust behavior requires a coarse-grained
pressureless current EFT; it is not a consequence of a winding integer alone.
Likewise, a fixed four-form supplies a vacuum-like stress tensor, but its flux
spacing and branch selection require additional data.

The value of the theorem is not that it produces a new homogeneous expansion
history, but that it gives a convention-explicit null test for when a
reconstructed Einstein residual may consistently be interpreted as a
pressureless protected current plus fixed four-form vacuum, while proving that
compact topology alone cannot select the observed amplitudes.  To become a
topology-evidence claim it must be supplemented by quotient transfer functions,
compact covariance kernels, likelihoods, and blind-injection tests.

\begin{thebibliography}{99}

\bibitem{Wald}
R.~M. Wald, \emph{General Relativity} (University of Chicago Press, Chicago, 1984).

\bibitem{HawkingEllis}
S.~W. Hawking and G.~F.~R. Ellis, \emph{The Large Scale Structure of Space-Time}
(Cambridge University Press, Cambridge, 1973).

\bibitem{Wolf}
J.~A. Wolf, \emph{Spaces of Constant Curvature}, 6th ed. (AMS Chelsea,
Providence, 2011).

\bibitem{Hatcher}
A. Hatcher, \emph{Algebraic Topology} (Cambridge University Press, Cambridge, 2002).

\bibitem{BrownFluid}
J.~D. Brown, Action functionals for relativistic perfect fluids,
Class. Quantum Grav. \textbf{10}, 1579--1606 (1993).


\bibitem{Aurilia}
A. Aurilia, H. Nicolai, and P.~K. Townsend,
Hidden constants: the theta parameter of QCD and the cosmological constant of
\(N=8\) supergravity, Nucl. Phys. B \textbf{176}, 509 (1980).

\bibitem{HenneauxTeitelboim}
M. Henneaux and C. Teitelboim,
The cosmological constant as a canonical variable,
Phys. Lett. B \textbf{143}, 415 (1984).

\bibitem{BrownTeitelboim}
J.~D. Brown and C. Teitelboim,
Neutralization of the cosmological constant by membrane creation,
Nucl. Phys. B \textbf{297}, 787 (1988).

\bibitem{BoussoPolchinski}
R. Bousso and J. Polchinski,
Quantization of four-form fluxes and dynamical neutralization of the cosmological constant,
JHEP \textbf{06}, 006 (2000).

\bibitem{Bardeen}
J.~M. Bardeen, Gauge-invariant cosmological perturbations,
Phys. Rev. D \textbf{22}, 1882 (1980).

\bibitem{KodamaSasaki}
H. Kodama and M. Sasaki, Cosmological perturbation theory,
Prog. Theor. Phys. Suppl. \textbf{78}, 1 (1984).

\bibitem{Mukhanov}
V. Mukhanov, H. Feldman, and R. Brandenberger,
Theory of cosmological perturbations,
Phys. Rep. \textbf{215}, 203 (1992).

\bibitem{MaBertschinger}
C.-P. Ma and E. Bertschinger,
Cosmological perturbation theory in the synchronous and conformal Newtonian gauges,
Astrophys. J. \textbf{455}, 7 (1995).

\bibitem{Cornish}
N.~J. Cornish, D.~N. Spergel, G.~D. Starkman, and E. Komatsu,
Constraining the topology of the Universe,
Phys. Rev. Lett. \textbf{92}, 201302 (2004).

\bibitem{PlanckTopology}
P.~A.~R. Ade \emph{et al.} (Planck Collaboration),
Planck 2015 results. XVIII. Background geometry and topology,
Astron. Astrophys. \textbf{594}, A18 (2016), arXiv:1502.01593.

\bibitem{Planck2018}
N. Aghanim \emph{et al.} (Planck Collaboration),
Planck 2018 results. VI. Cosmological parameters,
Astron. Astrophys. \textbf{641}, A6 (2020), arXiv:1807.06209.

\bibitem{SahaLens}
S. Saha \emph{et al.} (COMPACT Collaboration),
Cosmic topology. Part Ic. Limits on lens spaces from circle searches,
JCAP \textbf{01}, 004 (2025), arXiv:2409.02226.

\bibitem{DESIDR2}
M.~Abdul Karim \emph{et al.} (DESI Collaboration),
DESI DR2 results II: Measurements of baryon acoustic oscillations and cosmological constraints,
Phys. Rev. D \textbf{112}, 083515 (2025), arXiv:2503.14738.

\end{thebibliography}

\end{document}
